\section{Валидность и границы обобщения}
\label{sec:heldout-limits}
Первый и расширенный факторные эксперименты используют две модели, каждая с одним уровнем рассуждения и одной версией клиента. Выборки и число повторов различаются. Это расширяет область наблюдений по сравнению с одной моделью, но не позволяет выделить эффект выбора модели или обобщить результат на семейства моделей. Разработческие повторы остаются отдельной проверкой осуществимости. Публичность задач не позволяет исключить их присутствие в обучении; сходные задачи могут нарушать допущение независимости. Другая модель, язык, распределение задач или версия сервиса могут изменить направление эффекта. Версия клиента и даты запусков не фиксируют веса модели поставщика.

Лимит времени и административное продолжение задают ещё одну границу: неудачные генерации остаются пропусками, а изменение правила исполнения раскрыто. Поэтому вывод по доступным парам относится к наблюдаемой оцениваемой части; границы по всем назначениям показывают крайние варианты ненаблюдаемых исходов, но не устраняют смещение отбора. Превышение времени может отражать вычислительную сложность, а не случайный сетевой сбой, так что отсутствие обнаруженной разности не доказывает эквивалентность. Без внешне обоснованной границы допустимой потери и совместного критерия качества и стоимости решения о неухудшении или денежном превосходстве не принимаются.

Воздействия относятся к явным инструкциям. Названия ролей не удостоверяют наличие независимых агентов, а границы вызовов также раскрывают промежуточные ответы. Информация о задаче и операции этапов контролируются, но общий жёсткий лимит выходных токенов отсутствует. Ресурсные сравнения включают индуцированную длину рассуждений, повторную передачу контекста и накладные расходы клиента. Условная денежная оценка зависит от кеша; чувствительность без кеша убирает скидку, но не делает исполнение через подписку тождественным API-развёртыванию. Исследовательские вычисления и оценивание исключены; совокупная стоимость владения не оценивается.

Нативные тесты повышают проверяемость, сохраняя ограничения бенчмарка. Эталонный и неправильный контроли выявляют непригодную среду или нечувствительную отрицательную проверку, но их прохождение не доказывает соответствие каждого исходного утверждения естественно-языковой задаче. Отдельный аудит этапа разработки фиксирует расхождения инструкций и тестов как аннотации; исходные промпты, тесты и оценки остаются частью описания протокола. Эти аннотации не служат исключениями после получения результатов или новым эталоном. Семь задач с неработающим эталоном в первой серии mini и 15 таких задач в серии Luna сохраняются в соответствующем учёте назначений и ресурсов.

Сравнение INoT относится к раскрытой промптовой адаптации с явным разрешением неоднозначности источника. Её отдельная серия сохраняет различия времени и кеша. Завершённое расширение Luna является отдельным основным исследованием: 15\,000 назначений, 14\,961 завершённая программа и 985 задач, допускаемых контрольными проверками; оно повторно использует первые 200 идентификаторов и поэтому не является новой отложенной выборкой или репликацией при неизменных весах. Self-Collaboration (SCC) сравнивается по изменённому префиксу из 2\,000 блоков (6\,000 назначений трёх методов); генерация продолжается, нативные исходы ещё не просмотрены. Восемь ранее отправленных прерванных попыток остаются неизвестными на паузе и не повторяются. Авторские реализации INoT, Agentless и SWE-agent также не оценивались. Последние системы добавляют поиск, инструменты и взаимодействие с репозиторием за пределами факторного вопроса с фиксированным контекстом. Проба SWE-bench на одном issue даёт настоящие патчи и тестовые исходы, но полная матрица десяти назначений на одной задаче не позволяет сообщать оценку репозиторного бенчмарка или рейтинг агентов. Для таких широких выводов нужны нативное сравнение на нескольких репозиториях и репликация на нескольких моделях при сопоставимом бюджете.
